\SourceLeaf{393}%
Geraden. In diesem Fall lässt sich aber leicht zeigen, dass der gemeinsame Mittelpunkt notwendig auch der Mittelpunkt der Ellipse selbst ist.

2) Verlängert man zwei gegenüberliegende Seiten des Vierecks bis zu ihrem Schnittpunkt und ebenso die beiden anderen, so kann man auch die Strecke zwischen diesen beiden Schnittpunkten als verschwindende Ellipse auffassen, die die vier Seiten des Vierecks berührt. Ihr Mittelpunkt muss also auf derselben Geraden liegen, die die Mittelpunkte der beiden Diagonalen verbindet. Diese allgemeine Eigenschaft jedes Vierecks ist meines Wissens bisher noch nicht bemerkt worden. Ich werde sie weiter unten auf einfache Weise direkt beweisen.

Um die Rechnungen noch weiter abzukürzen, nehme ich nun an, dass diese Gerade selbst als Abszissenachse gewählt wurde und folglich \(\varphi=0\) gilt. Der Ursprung der Abszissen bleibt wie zuvor beliebig. Gerade die Wahl \(\varphi=0\) macht eine der vier Grundgleichungen entbehrlich. Zur Bestimmung der vier unbekannten Größen \(t,u,r,\psi\) haben wir also zum einen die drei Gleichungen
\[
\begin{aligned}
2aa+t-4ar\cos A+rr\cos2A-u\cos2(A-\psi)&=0\\
2a'a'+t-4a'r\cos A'+rr\cos2A'-u\cos2(A'-\psi)&=0\\
2a''a''+t-4a''r\cos A''+rr\cos2A''-u\cos2(A''-\psi)&=0
\end{aligned}
\]
und zum anderen die Bedingung, dass der Flächeninhalt der Ellipse, dem das Produkt \(\alpha\beta\) offenbar proportional ist. und folglich auch \(4\alpha\alpha\beta\beta\) oder \((rr-t)^2-uu\) maximal sein soll.

Setzt man der Kürze halber \(rr-t=\theta\) und
\[
\begin{aligned}
b&=2(a-r\cos A)^2\\
b'&=2(a'-r\cos A')^2\\
b''&=2(a''-r\cos A'')^2
\end{aligned}
\]
so werden die obigen Gleichungen zu
\[
\begin{aligned}
\theta+u\cos2(A-\psi)&=b\\
\theta+u\cos2(A'-\psi)&=b'\\
\theta+u\cos2(A''-\psi)&=b''
\end{aligned}
\]
